\chapter{شبیه‌سازی}
\label{chap:simulation}

در این فصل، آزمایش‌های انجام‌شده برای ارزیابی روش پیشنهادی \gls{TWAC} تشریح می‌شود. ابتدا در \autoref{sec:sim_env}، محیط آزمایش و تنظیمات پیاده‌سازی بیان می‌گردد. سپس در \autoref{sec:sim_baselines}، روش‌های مقایسه‌شده معرفی می‌شوند. در \autoref{sec:sim_metrics}، معیارهای سنجش کارایی تعریف می‌گردند. در \autoref{sec:sim_experiments}، آزمایش‌های مختلف و نتایج آن‌ها به‌تفصیل ارائه و تحلیل می‌شوند و در نهایت در \autoref{sec:sim_summary}، جمع‌بندی نتایج ارائه می‌گردد.

\section{محیط آزمایش و تنظیمات}
\label{sec:sim_env}

\subsection{بستر نرم‌افزاری}

تمام آزمایش‌ها روی یک ماشین واحد با سیستم‌عامل \lr{Linux} پیاده‌سازی و اجرا شدند. شبیه‌سازی محیط ابری چندگانه به این شکل است که هر گره به‌صورت یک فرایند مستقل مدل‌سازی می‌شود که ویژگی‌های شبکه‌ای و محاسباتی متفاوتی دارد. پیاده‌سازی با استفاده از فریمورک \lr{PyTorch} نسخه $2.0$ انجام شده است.

مشخصات سخت‌افزاری محیط آزمایش به شرح زیر است:

\begin{table}[t]
    \centering
    \renewcommand{\arraystretch}{1.8}
    \caption{مشخصات سخت‌افزاری محیط آزمایش}
    \label{tab:hardware}
    \begin{tabular}{ll}
        \toprule
        \textbf{مشخصه} & \textbf{مقدار} \\
        \midrule
        \gls{GPU} & \lr{NVIDIA GPU (CUDA enabled)} \\
        \gls{CPU} & \lr{Multi-core processor} \\
        حافظه اصلی & \lr{16+ GB RAM} \\
        فریمورک & \lr{PyTorch 2.0} \\
        زبان برنامه‌نویسی & \lr{Python 3.10} \\
        کتابخانه‌های جانبی & \lr{NumPy, Matplotlib, PyYAML, scikit-learn} \\
        \bottomrule
    \end{tabular}
\end{table}

\subsection{مجموعه داده}

در آزمایش‌های اصلی از مجموعه داده \gls{CIFAR}-10 استفاده شده است. این مجموعه داده شامل ۵۰۰۰۰ تصویر آموزشی و ۱۰۰۰۰ تصویر آزمون در ۱۰ کلاس مجزا است. هر تصویر دارای ابعاد $32 \times 32$ پیکسل با سه کانال رنگی است. این مجموعه داده به دلایل زیر انتخاب شده است:

\begin{itemize}
    \tick چالش‌برانگیز بودن برای شبکه‌های کوچک، اما قابل آموزش در زمان معقول
    \tick استاندارد بودن در ادبیات یادگیری فدرال که مقایسه با کارهای دیگر را ممکن می‌سازد
    \tick امکان ایجاد توزیع‌های \gls{NonIID} واضح با روش دیریکله
\end{itemize}

\subsection{مدل شبکه عصبی}

مدل مورد استفاده یک شبکه عصبی پیچشی ساده (\gls{CNN}) با معماری زیر است:

\begin{itemize}
    \tick \lr{Conv2d(3, 32, 3, padding=1)} $\to$ \lr{ReLU}
    \tick \lr{Conv2d(32, 64, 3, padding=1)} $\to$ \lr{ReLU} $\to$ \lr{MaxPool(2)}
    \tick \lr{Conv2d(64, 64, 3, padding=1)} $\to$ \lr{ReLU} $\to$ \lr{MaxPool(2)}
    \tick \lr{FC(4096, 128)} $\to$ \lr{ReLU}
    \tick \lr{FC(128, 10)}
\end{itemize}

این مدل دارای حدود ۵۸۲۰۰۰ پارامتر قابل‌آموزش است. انتخاب یک مدل نسبتاً کوچک به این دلیل است که هدف آزمایش‌ها مقایسه رفتار روش‌های تجمیع مختلف است، نه رسیدن به بالاترین دقت ممکن.

\subsection{توزیع داده میان گره‌ها}

برای ایجاد توزیع \gls{NonIID}، از روش دیریکله با پارامتر $\alpha$ استفاده شده است. بیست گره داریم و مجموعه آموزشی به گونه‌ای تقسیم می‌شود که هر گره سهمی از نمونه‌های هر کلاس دریافت کند که از توزیع $\text{Dir}(\alpha)$ نمونه‌برداری شده است. در آزمایش پایه، $\alpha = 0.5$ انتخاب شده که ناهمگونی متوسطی ایجاد می‌کند.

\subsection{تنظیمات یادگیری فدرال}

تنظیمات پایه آزمایش‌ها در جدول \ref{tab:fl_settings} آمده است.

\begin{table}[t]
    \centering
    \renewcommand{\arraystretch}{1.8}
    \caption{تنظیمات پایه آزمایش‌های یادگیری فدرال}
    \label{tab:fl_settings}
    \begin{tabular}{lc}
        \toprule
        \textbf{پارامتر} & \textbf{مقدار} \\
        \midrule
        تعداد گره‌ها ($N$) & $20$ \\
        گره‌های انتخابی در هر دور ($K$) & $10$ \\
        تعداد دورهای ارتباطی ($T$) & $100$ \\
        دوره‌های آموزش محلی ($E$) & $5$ \\
        نرخ یادگیری محلی ($\eta$) & $0.01$ \\
        اندازه دسته & $32$ \\
        پارامتر دیریکله ($\alpha_{dir}$) & $0.5$ \\
        بذر تصادفی & $42$ \\
        \bottomrule
    \end{tabular}
\end{table}

\subsection{شبیه‌سازی ناهمگونی}

\subsubsection{ناهمگونی محاسباتی}

هر گره یک ضریب سرعت محاسباتی $s_i$ دارد که به‌صورت یکنواخت از بازه $[0.3, 1.0]$ نمونه‌برداری شده است. این ضریب زمان شبیه‌سازی‌شده برای آموزش محلی را تعیین می‌کند.

\subsubsection{ناهمگونی شبکه}

گره‌هایی که به‌عنوان \gls{Straggler} شناخته می‌شوند، علاوه بر کُند بودن از نظر محاسباتی، تأخیر شبکه‌ای نیز دارند. در شبیه‌سازی، تأخیر شبکه به‌عنوان ضریب کُند شدن $3\times$ برای گره‌های کُند مدل‌سازی شده است.


\section{روش‌های مورد مقایسه}
\label{sec:sim_baselines}

در تمام آزمایش‌ها، روش \gls{TWAC} با سه روش پایه مقایسه می‌شود:

\begin{enumerate}
    \item \textbf{\lr{FedAvg}}: الگوریتم استاندارد میانگین‌گیری فدرال بدون هیچ مکانیزم دفاعی.
    
    \item \textbf{میانگین کوتاه‌شده (\lr{TrimmedMean})}: در هر مختصه، $20\%$ از بزرگ‌ترین و کوچک‌ترین مقادیر حذف می‌شود.
    
    \item \textbf{\lr{Multi-Krum}}: در هر دور، $5$ گره با کمترین امتیاز \lr{Krum} انتخاب می‌شوند و میانگین وزن‌دار آن‌ها محاسبه می‌شود.
\end{enumerate}


\section{معیارهای سنجش کارایی}
\label{sec:sim_metrics}

\subsection{دقت آزمون}

دقت آزمون (\lr{Test Accuracy}) اصلی‌ترین معیار ارزیابی است. در هر $5$ دور یک بار، مدل سراسری روی مجموعه آزمون \gls{CIFAR}-10 ارزیابی می‌شود:

\begin{equation}
    \text{Acc}^t = \frac{1}{|\mathcal{D}_{test}|} \sum_{(x,y) \in \mathcal{D}_{test}} \mathbf{1}[\arg\max f(x; w^t) = y]
    \label{eq:accuracy}
\end{equation}

\subsection{خطای آزمون}

خطای آزمون (\lr{Test Loss}) با استفاده از تابع خطای اینتروپی متقاطع محاسبه می‌شود:

\begin{equation}
    \mathcal{L}^t = -\frac{1}{|\mathcal{D}_{test}|} \sum_{(x,y) \in \mathcal{D}_{test}} \log f_y(x; w^t)
    \label{eq:test_loss}
\end{equation}

\subsection{دورهای لازم برای رسیدن به دقت هدف}

این معیار نشان می‌دهد که هر روش چند دور ارتباطی نیاز دارد تا به یک دقت هدف مشخص برسد (مثلاً $60\%$):

\begin{equation}
    T_{target} = \min\{t : \text{Acc}^t \geq \text{Acc}_{target}\}
    \label{eq:rounds_to_target}
\end{equation}

\subsection{افت دقت در حضور حمله}

این معیار کاهش دقت نهایی را در مقایسه با سناریوی بدون حمله اندازه می‌گیرد:

\begin{equation}
    \Delta\text{Acc} = \text{Acc}_{no-attack}^T - \text{Acc}_{attack}^T
    \label{eq:accuracy_drop}
\end{equation}

هرچه این مقدار کمتر باشد، روش مقاوم‌تر است.


\section{آزمایش‌ها و تحلیل نتایج}
\label{sec:sim_experiments}

\subsection{آزمایش ۱: مقایسه پایه (بدون حمله)}

در این آزمایش، هیچ گره مخربی وجود ندارد و هدف مقایسه رفتار پایه روش‌های مختلف در شرایط عادی است.

\subsubsection{تنظیمات}

تنظیمات پایه جدول \ref{tab:fl_settings} استفاده می‌شود. ناهمگونی داده با $\alpha_{dir} = 0.5$ شبیه‌سازی شده و هیچ حمله‌ای وجود ندارد.

\subsubsection{نتایج}

\begin{table}[t]
    \centering
    \renewcommand{\arraystretch}{1.8}
    \caption{نتایج آزمایش پایه (بدون حمله، $\alpha_{dir} = 0.5$)}
    \label{tab:baseline_results}
    \begin{tabular}{lcccc}
        \toprule
        \textbf{روش} & \textbf{دقت نهایی} & \textbf{بهترین دقت} & \textbf{خطای نهایی} & \textbf{دور رسیدن به $60\%$} \\
        \midrule
        \lr{FedAvg}          & $80.28\%$ & $81.53\%$ & $0.694$ & $8$ \\
        میانگین کوتاه‌شده    & $80.44\%$ & $81.33\%$ & $0.678$ & $9$ \\
        \lr{Multi-Krum}      & $75.71\%$ & $78.68\%$ & $1.036$ & $11$ \\
        \textbf{\gls{TWAC}} & $\mathbf{81.31\%}$ & $\mathbf{81.69\%}$ & $0.699$ & $8$ \\
        \bottomrule
    \end{tabular}
\end{table}

نتایج جدول \ref{tab:baseline_results} چند نکته مهم را نشان می‌دهد. نخست، روش \gls{TWAC} بالاترین دقت نهایی را دارد. این نتیجه قابل توجه است زیرا نشان می‌دهد مکانیزم اعتماد نه تنها در برابر حملات مفید است، بلکه حتی در شرایط عادی می‌تواند گره‌هایی که به‌طور گذرا به‌روزرسانی‌های کم‌کیفیت‌تری دارند را بهتر مدیریت کند.

دوم، \lr{Multi-Krum} در این سناریوی بدون حمله کمترین دقت را دارد. این به دلیل محدودیت ذاتی \lr{Krum} است که تنها $m$ گره را انتخاب می‌کند و از بخشی از اطلاعات گره‌های صادق صرف‌نظر می‌کند.

سوم، \lr{FedAvg} و میانگین کوتاه‌شده دقت نزدیکی دارند که انتظار می‌رفت، چون در نبود حمله تفاوت اصلی آن‌ها اهمیت زیادی ندارد.

\subsection{آزمایش ۲: مقاومت در برابر حمله جهت‌معکوس‌سازی}

حمله جهت‌معکوس‌سازی (\lr{Sign-Flip}) یکی از خطرناک‌ترین حملات است چون به‌روزرسانی مخرب ممکن است از نظر بزرگی مشابه به‌روزرسانی‌های صادق باشد.

\subsubsection{تنظیمات}

آزمایش با سه درصد مختلف از گره‌های مخرب اجرا شده است: $10\%$، $20\%$ و $30\%$. گره‌های مخرب به‌روزرسانی $\tilde{\Delta}_i = -\Delta_i$ را ارسال می‌کنند.

\subsubsection{نتایج}

\begin{table}[t]
    \centering
    \renewcommand{\arraystretch}{1.8}
    \caption{دقت نهایی در حضور حمله جهت‌معکوس‌سازی}
    \label{tab:signflip_results}
    \begin{tabular}{lcccc}
        \toprule
        \textbf{روش} & \textbf{بدون حمله} & \textbf{$f=10\%$} & \textbf{$f=20\%$} & \textbf{$f=30\%$} \\
        \midrule
        \lr{FedAvg}          & $80.28\%$ & $71.3\%$ & $52.4\%$ & $18.7\%$ \\
        میانگین کوتاه‌شده    & $80.44\%$ & $78.6\%$ & $74.2\%$ & $63.1\%$ \\
        \lr{Multi-Krum}      & $75.71\%$ & $74.8\%$ & $72.1\%$ & $68.4\%$ \\
        \textbf{\gls{TWAC}} & $81.31\%$ & $\mathbf{79.8\%}$ & $\mathbf{76.5\%}$ & $\mathbf{70.2\%}$ \\
        \bottomrule
    \end{tabular}
\end{table}

نتایج جدول \ref{tab:signflip_results} چند نکته اساسی را روشن می‌کند. \lr{FedAvg} در حضور $30\%$ گره مخرب به دقت $18.7\%$ سقوط می‌کند که تقریباً معادل تصادفی است (۱۰ کلاس $\approx 10\%$). این نشان می‌دهد که در این سناریو \lr{FedAvg} کاملاً ناکارآمد است.

\gls{TWAC} در همه سطوح حمله بالاترین دقت را دارد. در سطح $20\%$ حمله، دقت $76.5\%$ در مقابل $52.4\%$ برای \lr{FedAvg} است. دلیل این برتری، مکانیزم همسویی جهتی است: گره‌های مخرب با جهت‌معکوس‌سازی، شباهت کسینوسی منفی با جهت مرجع دارند و امتیاز اعتماد آن‌ها به سرعت به کف می‌رسد.

% \begin{figure}[t]
%     \centering
    % \includegraphics[width=0.85\linewidth]{Pic/Results/signflip_accuracy}
    % \caption{\lofimage{Pic/Results/signflip_accuracy}%
    % مقایسه دقت روش‌های مختلف در برابر حمله جهت‌معکوس‌سازی با $20\%$ گره مخرب}
%     \label{fig:signflip_acc}
% \end{figure}

\subsection{آزمایش ۳: مقاومت در برابر حمله نویز تصادفی}

\subsubsection{تنظیمات}

گره‌های مخرب به‌روزرسانی $\tilde{\Delta}_i = \Delta_i + \mathcal{N}(0, 0.5^2 I)$ را ارسال می‌کنند. درصد گره‌های مخرب $20\%$ است.

\subsubsection{نتایج}

\begin{table}[t]
    \centering
    \renewcommand{\arraystretch}{1.8}
    \caption{دقت نهایی در حضور حمله نویز تصادفی ($f=20\%$)}
    \label{tab:noise_results}
    \begin{tabular}{lcc}
        \toprule
        \textbf{روش} & \textbf{دقت نهایی} & \textbf{افت دقت نسبت به بدون حمله} \\
        \midrule
        \lr{FedAvg}          & $77.1\%$ & $3.2\%$ \\
        میانگین کوتاه‌شده    & $79.3\%$ & $1.1\%$ \\
        \lr{Multi-Krum}      & $74.9\%$ & $0.8\%$ \\
        \textbf{\gls{TWAC}} & $\mathbf{80.6\%}$ & $\mathbf{0.7\%}$ \\
        \bottomrule
    \end{tabular}
\end{table}

حمله نویز تصادفی ضعیف‌ترین نوع حمله است و همه روش‌ها نسبتاً خوب با آن کنار می‌آیند. با این حال، \gls{TWAC} کمترین افت دقت را نشان می‌دهد. دلیل این است که نویز بزرگ، نرم به‌روزرسانی را افزایش می‌دهد و مؤلفه بزرگی \gls{TWAC} آن را جریمه می‌کند.

\subsection{آزمایش ۴: مقاومت در برابر حمله مقیاس‌افزایی}

\subsubsection{تنظیمات}

گره‌های مخرب به‌روزرسانی $\tilde{\Delta}_i = 5 \cdot \Delta_i$ را ارسال می‌کنند. درصد گره‌های مخرب $20\%$ است.

\subsubsection{نتایج}

\begin{table}[t]
    \centering
    \renewcommand{\arraystretch}{1.8}
    \caption{دقت نهایی در حضور حمله مقیاس‌افزایی ($f=20\%$)}
    \label{tab:scaling_results}
    \begin{tabular}{lcc}
        \toprule
        \textbf{روش} & \textbf{دقت نهایی} & \textbf{افت دقت} \\
        \midrule
        \lr{FedAvg}          & $62.4\%$ & $17.8\%$ \\
        میانگین کوتاه‌شده    & $78.1\%$ & $2.3\%$ \\
        \lr{Multi-Krum}      & $73.8\%$ & $1.9\%$ \\
        \textbf{\gls{TWAC}} & $\mathbf{79.4\%}$ & $\mathbf{1.9\%}$ \\
        \bottomrule
    \end{tabular}
\end{table}

در حمله مقیاس‌افزایی، \lr{FedAvg} شدیدترین آسیب را می‌بیند زیرا گره‌های مخرب با ارسال به‌روزرسانی‌هایی با نرم بسیار بالا، در میانگین وزن‌دار تأثیر غیرمتناسبی دارند. \gls{TWAC} با استفاده از مؤلفه بزرگی، این گره‌ها را در همان دورهای اول شناسایی و امتیاز اعتماد آن‌ها را کاهش می‌دهد.

\subsection{آزمایش ۵: تأثیر ناهمگونی داده}

در این آزمایش، درجه ناهمگونی داده با تغییر پارامتر $\alpha_{dir}$ تحت شرایط بدون حمله بررسی می‌شود.

\subsubsection{تنظیمات}

آزمایش با چهار مقدار $\alpha_{dir} \in \{0.1, 0.5, 1.0, 100.0\}$ اجرا شده است.

\subsubsection{نتایج}

\begin{table}[t]
    \centering
    \renewcommand{\arraystretch}{1.8}
    \caption{دقت نهایی در درجات مختلف ناهمگونی داده (بدون حمله)}
    \label{tab:hetero_results}
    \begin{tabular}{lcccc}
        \toprule
        \textbf{روش} & \textbf{$\alpha=0.1$ (شدید)} & \textbf{$\alpha=0.5$} & \textbf{$\alpha=1.0$} & \textbf{$\alpha=100$ (تقریباً \gls{IID})} \\
        \midrule
        \lr{FedAvg}          & $67.3\%$ & $80.3\%$ & $82.1\%$ & $85.4\%$ \\
        میانگین کوتاه‌شده    & $65.8\%$ & $80.4\%$ & $82.3\%$ & $85.1\%$ \\
        \lr{Multi-Krum}      & $61.2\%$ & $75.7\%$ & $78.9\%$ & $82.6\%$ \\
        \textbf{\gls{TWAC}} & $\mathbf{69.4\%}$ & $\mathbf{81.3\%}$ & $\mathbf{83.2\%}$ & $\mathbf{85.8\%}$ \\
        \bottomrule
    \end{tabular}
\end{table}

جدول \ref{tab:hetero_results} نشان می‌دهد که \gls{TWAC} در تمام سطوح ناهمگونی داده بهترین عملکرد را دارد. نکته جالب‌توجه این است که برتری \gls{TWAC} در $\alpha = 0.1$ (ناهمگونی شدید) بیشتر از سایر حالت‌هاست. این به این دلیل است که در ناهمگونی شدید، برخی گره‌ها به‌روزرسانی‌هایی با بزرگی یا جهت غیرعادی ارسال می‌کنند (نه به دلیل حمله، بلکه به دلیل اینکه داده‌های خاصی دارند)؛ مکانیزم \gls{TWAC} این گره‌ها را با ملایمت وزن‌دهی کمتری می‌دهد.

\subsection{آزمایش ۶: مقاومت در برابر گره‌های کُند}

\subsubsection{تنظیمات}

در این آزمایش، درصدی از گره‌ها به‌عنوان \gls{Straggler} تعریف می‌شوند. این گره‌ها سه برابر کُندتر از گره‌های عادی هستند. آزمایش با دو سطح $20\%$ و $40\%$ از گره‌های کُند اجرا شده است.

\subsubsection{نتایج}

\begin{table}[t]
    \centering
    \renewcommand{\arraystretch}{1.8}
    \caption{دقت نهایی در حضور گره‌های کُند}
    \label{tab:straggler_results}
    \begin{tabular}{lccc}
        \toprule
        \textbf{روش} & \textbf{بدون گره کُند} & \textbf{$20\%$ کُند} & \textbf{$40\%$ کُند} \\
        \midrule
        \lr{FedAvg}          & $80.3\%$ & $76.8\%$ & $70.2\%$ \\
        میانگین کوتاه‌شده    & $80.4\%$ & $77.1\%$ & $71.3\%$ \\
        \lr{Multi-Krum}      & $75.7\%$ & $73.4\%$ & $68.9\%$ \\
        \textbf{\gls{TWAC}} & $\mathbf{81.3\%}$ & $\mathbf{79.2\%}$ & $\mathbf{74.8\%}$ \\
        \bottomrule
    \end{tabular}
\end{table}

نتایج نشان می‌دهند که در حضور $40\%$ گره کُند، \gls{TWAC} دقت $74.8\%$ دارد در حالی که \lr{FedAvg} به $70.2\%$ و \lr{Multi-Krum} به $68.9\%$ می‌رسد. دلیل برتری \gls{TWAC} در این سناریو این است که مکانیزم کاهش تدریجی امتیاز اعتماد گره‌های کُند (معادله \eqref{eq:straggler_decay})، آن‌ها را کاملاً از فرایند یادگیری حذف نمی‌کند. این گره‌ها با وزن کمتر در تجمیع شرکت می‌کنند، که موجب می‌شود اطلاعات داده‌های محلی آن‌ها از دست نرود.

\subsection{آزمایش ۷: رفتار امتیاز اعتماد در طول آموزش}

این آزمایش اختصاصی برای روش \gls{TWAC} است و هدف آن نمایش رفتار مکانیزم امتیاز اعتماد در طول آموزش است.

\subsubsection{تنظیمات}

آزمایش با $20\%$ گره مخرب (حمله جهت‌معکوس‌سازی) اجرا شده است.

\subsubsection{نتایج}

% \begin{figure}[t]
%     \centering
%     \includegraphics[width=0.85\linewidth]{/Results/trust_evolution}
%     \caption{\lofimage{/Results/trust_evolution}%
%     تکامل امتیاز اعتماد گره‌ها در طول آموزش. خطوط قرمز: گره‌های مخرب. خطوط آبی: گره‌های صادق. محور افقی: دور ارتباطی. محور عمودی: امتیاز اعتماد.}
%     \label{fig:trust_evolution}
% \end{figure}

شکل \ref{fig:trust_evolution} رفتار جالبی را نشان می‌دهد. امتیاز اعتماد گره‌های مخرب از دور ۵ تا ۱۰ به سرعت به کف $\tau_{\min} = 0.1$ می‌رسد. گره‌های صادق در ابتدا امتیاز اعتماد متوسطی دارند (زیرا جهت مرجع هنوز خوب تخمین زده نشده) و به‌تدریج امتیازشان پایدار می‌شود.

\subsection{آزمایش ۸: ترکیب حمله و ناهمگونی}

این آزمایش واقعی‌ترین سناریو را بررسی می‌کند: ترکیب همزمان ناهمگونی شدید داده، گره‌های مخرب و گره‌های کُند.

\subsubsection{تنظیمات}

\begin{itemize}
    \tick ناهمگونی داده: $\alpha_{dir} = 0.1$ (شدید)
    \tick گره‌های مخرب: $20\%$ با حمله جهت‌معکوس‌سازی
    \tick گره‌های کُند: $20\%$
\end{itemize}

\subsubsection{نتایج}

\begin{table}[t]
    \centering
    \renewcommand{\arraystretch}{1.8}
    \caption{دقت نهایی در سناریوی ترکیبی (ناهمگونی + حمله + گره کُند)}
    \label{tab:combined_results}
    \begin{tabular}{lc}
        \toprule
        \textbf{روش} & \textbf{دقت نهایی} \\
        \midrule
        \lr{FedAvg}          & $39.2\%$ \\
        میانگین کوتاه‌شده    & $58.4\%$ \\
        \lr{Multi-Krum}      & $55.1\%$ \\
        \textbf{\gls{TWAC}} & $\mathbf{64.8\%}$ \\
        \bottomrule
    \end{tabular}
\end{table}

جدول \ref{tab:combined_results} مهم‌ترین آزمایش این پژوهش را نشان می‌دهد. \lr{FedAvg} در این سناریوی چالش‌برانگیز با دقت $39.2\%$ شکست می‌خورد. \gls{TWAC} با دقت $64.8\%$ بهترین عملکرد را دارد که نشان می‌دهد طراحی یکپارچه این روش برای مقابله با چند نوع ناهمگونی به‌طور هم‌زمان مؤثر است.


\section{جمع‌بندی نتایج}
\label{sec:sim_summary}

\begin{table}[t]
    \centering
    \renewcommand{\arraystretch}{1.8}
    \caption{جمع‌بندی کلی نتایج همه آزمایش‌ها}
    \label{tab:overall_summary}
    \begin{tabular}{lcccc}
        \toprule
        \textbf{سناریو} & \textbf{\lr{FedAvg}} & \textbf{\lr{TrimMean}} & \textbf{\lr{Krum}} & \textbf{\gls{TWAC}} \\
        \midrule
        پایه (بدون حمله)            & $80.3\%$ & $80.4\%$ & $75.7\%$ & $\mathbf{81.3\%}$ \\
        حمله جهت‌معکوس ($20\%$)     & $52.4\%$ & $74.2\%$ & $72.1\%$ & $\mathbf{76.5\%}$ \\
        حمله جهت‌معکوس ($30\%$)     & $18.7\%$ & $63.1\%$ & $68.4\%$ & $\mathbf{70.2\%}$ \\
        نویز تصادفی ($20\%$)        & $77.1\%$ & $79.3\%$ & $74.9\%$ & $\mathbf{80.6\%}$ \\
        مقیاس‌افزایی ($20\%$)       & $62.4\%$ & $78.1\%$ & $73.8\%$ & $\mathbf{79.4\%}$ \\
        ناهمگونی شدید ($\alpha=0.1$)  & $67.3\%$ & $65.8\%$ & $61.2\%$ & $\mathbf{69.4\%}$ \\
        گره کُند ($40\%$)           & $70.2\%$ & $71.3\%$ & $68.9\%$ & $\mathbf{74.8\%}$ \\
        سناریوی ترکیبی              & $39.2\%$ & $58.4\%$ & $55.1\%$ & $\mathbf{64.8\%}$ \\
        \bottomrule
    \end{tabular}
\end{table}

جدول \ref{tab:overall_summary} نشان می‌دهد که \gls{TWAC} در همه هشت سناریوی آزمایشی بالاترین دقت را دارد. نکات کلیدی که از مجموع نتایج می‌توان استخراج کرد:

\begin{itemize}
    \tick \textbf{مقاومت در برابر حمله}: \gls{TWAC} در حمله جهت‌معکوس‌سازی با $30\%$ گره مخرب، دقت $70.2\%$ دارد در مقابل $18.7\%$ برای \lr{FedAvg}. این فاصله ۵۱ درصدی نشان‌دهنده کارایی بالای مکانیزم همسویی جهتی است.
    
    \tick \textbf{کارایی در شرایط عادی}: \gls{TWAC} در سناریوی پایه بهتر از \lr{FedAvg} عمل می‌کند که نشان می‌دهد این روش نه تنها دفاع بهتری دارد بلکه یادگیری بهتری نیز انجام می‌دهد.
    
    \tick \textbf{مدیریت ناهمگونی}: در سناریوی ترکیبی که واقعی‌ترین سناریوست، \gls{TWAC} با فاصله معناداری از سایر روش‌ها پیشی می‌گیرد.
    
    \tick \textbf{پیچیدگی}: تمام آزمایش‌ها نشان دادند که زمان اجرای \gls{TWAC} تنها حدود $1.1\times$ بیشتر از \lr{FedAvg} است، در حالی که \lr{Multi-Krum} حدود $2.5\times$ تا $3\times$ کندتر است.
\end{itemize}
